\scanpage{92}\mspage{T-6}

\uline{Cas}\quad $j=0$\quad (i.e. $a_{2}=0$).

Forme réduite :
\[
 y^{2}=x^{3}+a_{4}x+a_{6}\qquad \handwritten{\Delta}=-a_{4\unsure{.|\phantom{.}}}^{3}
 \qquad \pi=a_{4}^{-1}dy
\]

Isomorphismes :
\[
 x\unsure{\mathord{.}|\circ}f=u^{2}x'+r\ ,\quad
 y\mathord{.}f=u^{3}y'
\]

avec :
\[
 u^{4}a'_{4}=a_{4}\ ,\quad u^{6}a'_{6}=a_{6}+ra_{4}+r^{3}
\]

Deux courbes $C$ et $C'$ de même $j$ sont isomorphes si et seulement si :
\[
 a_{4}/a'_{4}\in k^{\ast 4}\qquad\text{et}\quad
 (a_{4}/a'_{4})^{3/2}a'_{6}-a_{6}\quad\text{est de la forme }r^{3}+ra_{4}.
\]

Cette condition est en tout cas vérifiée sur une extension séparable $k'/k$ de degré $\handwritten{\leqslant}\,12$.

Les automorphismes de $C$ sont donnés par :
\[
\begin{aligned}
u&=\pm 1 & r^{3}+a_{4}r&\deleted{\unsure{+(a_{6}-a'_{6})|+a_{6}-a'_{6}}}=0\\
u&=\pm i & r^{3}+a_{4}r&+\deleted{\unsure{(|\phantom{(}}}\handwritten{2}a_{6}\deleted{\unsure{+a'_{6})|-a'_{6})|+a'_{6}|-a'_{6}}}=0\ .
\end{aligned}
\]

Ils forment un groupe d'ordre $12$ (sur la clôture séparable de $k$).

L'anneau d'endomorphismes est un ordre maximal d'un corps de quaternions ramifié en $3$ et $\handwritten{\infty}$ .

Courbe type :
\[
 y^{2}=x^{3}-x.
\]

Dans ce cas, les automorphismes sont donnés par :
\[
 u^{4}=1\ ,\quad r^{3}-r=0\quad\text{(i.e. }r\in\text{\uuline{F}}_{3}\text{)}.
\]
% archive sheet number: 91

